%==============================================================================
%  figures/observation_operator.tex
%  The observation operator: spatial interpolation to sensors (left) and
%  temporal averaging over intervals (right). Adapted from presentation.tex.
%  \input in the methodology observation-operator subsection.
%==============================================================================
\begin{figure}[tb]
  \centering
  \begin{subfigure}[b]{0.40\linewidth}
    \centering
    \begin{tikzpicture}[font=\tiny,scale=1.35,transform shape]
      \draw[uaCharcoal!22,line width=0.4pt] (0,0) grid (2,2);
      \foreach \x in {0,1,2}\foreach \y in {0,1,2}
        \fill[uaCharcoal!45] (\x,\y) circle (1.5pt);
      \coordinate (s) at (0.6,0.55);
      \draw[uaOrange!70,line width=0.6pt,densely dashed] (s)--(0,0);
      \draw[uaOrange!70,line width=0.6pt,densely dashed] (s)--(1,0);
      \draw[uaOrange!70,line width=0.6pt,densely dashed] (s)--(0,1);
      \draw[uaOrange!70,line width=0.6pt,densely dashed] (s)--(1,1);
      \fill[uaOrange] (s) circle (2.4pt);
      \node[text=uaOrange,anchor=south west,inner sep=1pt] at (s) {sensor};
      \node[text=uaGrey,anchor=north] at (1,-0.2) {grid nodes};
    \end{tikzpicture}
    \caption{Spatial: interpolation from grid nodes to $\Nd$ sensors.}
    \label{fig:obsop:spatial}
  \end{subfigure}
  \hfill
  \begin{subfigure}[b]{0.56\linewidth}
    \centering
    \begin{tikzpicture}[font=\tiny,>=latex,scale=1.4,transform shape]
      % bin edges
      \foreach \x in {0.5,1.0,...,4.0}
        \draw[uaCharcoal!25,line width=0.5pt,densely dashed] (\x,0.55) -- (\x,2.0);
      % axes
      \draw[->,uaCharcoal!70,line width=0.7pt] (0,0.55) -- (4.75,0.55);
      \node[text=uaGrey,anchor=north] at (0,0.5) {0};
      \node[text=uaGrey,anchor=north] at (4.5,0.5) {180 s};
      % bin bracket over first bin
      \draw[uaGrey,line width=0.6pt] (0,2.12) -- (0,2.22) -- (0.5,2.22) -- (0.5,2.12);
      \node[text=uaGrey,anchor=south] at (0.25,2.24) {10 s};
      % sensor signal (fine cadence)
      \draw[uaTeal,line width=0.9pt] plot[smooth] coordinates
        {(0,1.30) (0.25,1.05) (0.5,0.85) (0.75,1.00) (1.0,1.35)
         (1.25,1.55) (1.5,1.45) (1.75,1.20) (2.0,1.10) (2.25,1.25)
         (2.5,1.50) (2.75,1.65) (3.0,1.55) (3.25,1.30) (3.5,1.05)
         (3.75,0.90) (4.0,1.00) (4.25,1.20) (4.5,1.35)};
      % per-bin means -> observations
      \foreach \xa/\m in {0.0/1.07, 0.5/1.07, 1.0/1.45, 1.5/1.25,
                          2.0/1.28, 2.5/1.57, 3.0/1.30, 3.5/0.98, 4.0/1.18}{
        \draw[uaOrange,line width=1.4pt] (\xa+0.05,\m) -- (\xa+0.45,\m);
        \fill[uaOrange] (\xa+0.25,\m) circle (1.6pt);}
      % legend
      \draw[uaTeal,line width=0.9pt] (0.1,0.15) -- (0.45,0.15);
      \node[text=uaGrey,anchor=west] at (0.5,0.15) {sensor signal};
      \draw[uaOrange,line width=1.4pt] (2.3,0.15) -- (2.65,0.15);
      \fill[uaOrange] (2.475,0.15) circle (1.6pt);
      \node[text=uaGrey,anchor=west] at (2.7,0.15) {bin mean = one obs.};
    \end{tikzpicture}
    \caption{Temporal: each sensor signal is averaged over fixed intervals; one
      bin mean becomes one observation.}
    \label{fig:obsop:temporal}
  \end{subfigure}
  \caption{The observation operator \Hobs\ maps the model field to the data
    vector \obs\ in two steps: spatial interpolation to street-level sensors
    (a) and temporal averaging over intervals (b). Averaging suppresses the
    turbulent fluctuations so that \esmda\ targets the mean forcing.
    \todo{state $\Nd$, sensor count, interval length, and window relation.}}
  \label{fig:obs_operator}
\end{figure}
